\section{Related Literature}\label{sec:literature}

Spatial asset pricing asks how local interaction modifies factor-based pricing
once the researcher supplies an interaction matrix $W$.
Across geographic and financial applications, $W$ organizes local dependence
before its strength is estimated
\citep{fernandez_spatial_2011,kou_asset_2018,ge_news-implied_2023}.
In the spatial CAPM and spatial APT of \citet{kou_asset_2018}, assets occupy
point-valued locations and inverse geographic distance supplies the weights.
The resulting spatial multiplier accumulates direct and higher-order
interactions, but the asset representation and the rule that converts it into
$W$ remain exogenous to the pricing model.

Linear-quadratic network games address a different part of this problem.
Costly local complementarity yields best responses summarized by a network
resolvent, providing an economic route from individual objectives to aggregate
propagation conditional on a given graph \citep{ballester_networks_2006}.
The adjustment model in this paper applies that logic one layer below returns:
firms trade off departures from their stand-alone factor exposures against
misalignment among peer-adjusted exposures, and the spatial coefficient records
the relative intensity of that adjustment.
This interpretation explains why interaction may arise, but it leaves open
which firms should be peers and how their weights should be determined.

The news-implied-network literature supplies one influential answer by replacing
geographic location with an observed informational edge.
Beginning with economic linkages inferred from news and their relation to return
predictability, this line uses sentence-level co-mentions to construct firm
networks that trace contagion and aggregate risk
\citep{scherbina_economic_2013,schwenkler_network_2020}.
\citet{ge_news-implied_2023} places such a co-mention matrix inside a spatial
factor model, and subsequent work studies whether auxiliary network information
improves covariance estimation \citep{ge_network_guided_covariance_2026}.
These studies expand the meaning of economic proximity beyond literal distance,
but the observed or count-weighted graph remains the empirical primitive that
supplies $W$.

A broader text-finance literature usually maps information into a first-order
object such as an attention measure, a textual factor, or an embedding-based
input to return or stochastic-discount-factor estimation
\citep{ben_rephael_information_consumption_2019,
cong_textual_factors_2024,wang2025newsnet}.
Related learned-network methods map connectivity directly into factor exposures
or priced network factors \citep{son2022graph,uddin2024network}.
Together, these approaches establish that text and learned representations carry
economically relevant predictive and pricing information.
The distinct question here is not whether text predicts a scalar outcome, but
whether the full cross-section of within-firm information can construct the peer
field along which exposures adjust.

That question changes the primitive representation of a firm.
A diversified firm need not occupy one geographic or characteristic point;
its articles can instead be represented by a probability distribution $C_i$
over positions in a maintained information space.
This representation preserves dispersion, multimodality, and other differences
that a centroid or single embedding suppresses.
When every $C_i$ degenerates to a point mass, Wasserstein distance reduces to
the underlying point distance, so point-based spatial logic is nested at the
level of the representation.
For genuinely distribution-valued firms, however, the graph becomes an output
of the distributional geometry rather than its starting point.

Quadratic Wasserstein transport gives that geometry economic content beyond a
generic similarity metric.
Its primal problem finds the least aggregate squared displacement required to
reallocate one distribution into another.
Unlike the $W_1$ Kantorovich--Rubinstein case, the quadratic $W_2$ dual is not a
single Lipschitz price schedule, so it does not carry the same direct
spatial-arbitrage interpretation.
In this application, the ground cost is displacement in a maintained embedding
space rather than a monetary shipping cost, so the construction neither prices
literal transport nor tests for semantic arbitrage.
It measures least semantic displacement conditional on the chosen
representation and ground metric.

Pairwise transport nevertheless stops short of an interaction field.
A distance answers how much reallocation separates two firms, and a transport
plan establishes article-level correspondence, but neither determines how
several candidate firms should jointly represent one target.
Turning pairwise distances directly into weights would require an additional
kernel, bandwidth, or nearest-neighbour rule.
The central abstraction of this paper is instead
\emph{target-anchored Wasserstein barycentric reconstruction}.
For each target firm, quadratic transport first aligns its articles separately
with those of every candidate firm.
Holding those target-specific correspondences fixed, a convex simplex step
selects the nonnegative unit-sum distributional spanning weights that best
reconstruct the target cloud from its aligned peers.
The fitted rows form the \emph{barycentric interaction field} $W^\flat$.

The closest geometric reference is the unrestricted Wasserstein barycentre,
which selects a free centre distribution for several input laws
\citep{agueh_carlier_barycenters_2011}.
The present construction instead holds the target and pairwise alignments fixed
and estimates target-specific distributional spanning weights.
Section~\ref{sec:interaction-operator} states this boundary formally.
The distinction is economically consequential: a large $W^\flat_{ij}$ records
firm $j$'s conditional usefulness in reconstructing target $i$ from the
investable universe, not merely small pairwise distance.
Because that usefulness is target anchored, the barycentric interaction field
can be directed even though pairwise Wasserstein distance is symmetric.
Its direction records reconstruction relevance, not causal influence.

At neighboring levels of aggregation, related studies use distribution-valued
firm characteristics for different finance questions.
\citet{gawronsky_continuous_2026} studies pairwise covariance envelopes implied
by distances between characteristic laws, whereas
\citet{gawronsky_distance_mpt_2026} studies portfolio-risk bounds and allocation
from distributional structure.
The present paper occupies the intermediate, multi-firm level: it turns
target-specific correspondences into a cross-sectional field and studies the
propagation of stand-alone exposures through that field.
Its field construction and adjustment model are stated independently of the
pairwise and portfolio results, which locate the contribution without supplying
a premise for it.

Constructing the field from text also preserves the conditioning requirement of
spatial inference.
Classical spatial-autoregressive methods condition on a known $W$, whereas
estimating $W$ from the same outcomes used in the spatial lag creates mechanical
reflection \citep{anselin_1988,lesage_pace_2009,kelejian_prucha_2010}.
Freezing every text-derived field before the return-evaluation window removes
that same-sample feedback.
It does not identify causal peer effects when omitted industries, technologies,
attention, or selection jointly influence text and returns.

The possibility of several admissible fields produces the final change in the
literature's empirical question.
Estimating candidate matrices separately asks ``which $W$ wins?'' but cannot
distinguish redundant descriptions of the same relations from distinct channels
of exposure adjustment.
The multi-field quadratic model instead places the barycentric interaction field
beside a conventional news-link field and assigns each its own coefficient.
Separate-network specifications become boundary cases of the joint model, and
the estimand becomes how adjustment divides across channels, including the case
in which one field absorbs the other.
With the source of the field, the exposure-adjustment mechanism, and the return
bridge kept distinct, the next section introduces the stand-alone and
peer-adjusted exposures linked by that mechanism.
